\chapter{Mathematica Interface}

FEMM can be driven from Mathematica using the {\tt MathFEMM} package, which
controls a FEMM process over a TCP/IP socket.  (Earlier versions used
Mathematica's MathLink API; the current package uses the same TCP interface as
all of the other language bindings --- see the {\em TCP/IP Interface} chapter.)
The package takes care of launching FEMM, establishing the connection, sending
commands, and parsing the results, so that FEMM is controlled directly through a
set of Mathematica functions.

To use the package, load it and open a connection:

\begin{verbatim}
Get["mathfemm.m"]
OpenFEMM[]
\end{verbatim}

\verb+OpenFEMM[]+ connects to a FEMM process, launching one if necessary.  By
default it starts the interactive shell ({\tt femm\_gui}) and uses port 48015.
Pass \verb+OpenFEMM[1]+ to launch the headless engine ({\tt femm\_headless})
instead, with no window, or \verb+OpenFEMM[bHide, port]+ to choose both the
visibility and the port.

If a FEMM process is already listening on that port, {\tt OpenFEMM} attaches to
it rather than starting a second one.  That is what allows a notebook to be
re-run after an aborted script without leaving a trail of orphaned FEMM
windows, and it means the {\tt bHide} argument has no effect when an instance
is already running.

Once connected, FEMM is controlled through the Mathematica functions provided by
the package, which mirror the Python command set documented in the scripting
chapters.  For example, the magnetics preprocessor commands appear as
\verb+MIAddNode[x,y]+, \verb+MIAddSegment[x1,y1,x2,y2]+,
\verb+MIAddBlockLabel[x,y]+, and so on.  Each function sends the corresponding
command to FEMM's Python interpreter and returns the result as a Mathematica
expression --- a number, a string, or a list, as appropriate --- so values can
be used directly in further Mathematica computation.

When finished, close the connection with:

\begin{verbatim}
CloseFEMM[]
\end{verbatim}

\chapter{TCP/IP Interface}

Beyond the supplied language bindings, FEMM can be controlled by any program
that can open a TCP/IP socket.  A running FEMM process exposes its Python
interpreter over a socket, so an external client can send commands and read back
results programmatically.  (This replaces the ActiveX/COM automation server used
by older, Windows-only versions of FEMM.)

A FEMM process listens for a connection on a TCP port; the default is 48015.
The headless engine is started for this purpose with, for example:

\begin{verbatim}
femm_headless --port 48015
\end{verbatim}

The interactive shell {\tt femm\_gui} also serves the same protocol while
displaying its window.  The connection is single-client: one program controls a
given FEMM process at a time, and a second program that connects while the
first is still attached is refused with an error rather than queued.

The process exits when its client disconnects, so that a binding which is
killed does not leave an orphaned FEMM behind.  A consequence worth knowing:
a script that wants to leave the window up for inspection when it finishes
must keep its connection open, rather than closing it or calling
{\tt closefemm}.

The protocol is a simple line-based text exchange.  The client sends a command
--- any statement or expression accepted by FEMM's Python interpreter --- as a
single newline-terminated line, and FEMM replies with exactly one
newline-terminated line:

\begin{itemize}
\item If the command is an expression, the reply is its value rendered as text.
      Real numbers are reported in full precision; a complex number is returned
      in the form \verb|a+I*b|; and a tuple of results is returned with the
      individual items separated by tab characters.
\item A statement that produces no value (an assignment, a drawing command, and
      the like) replies with an empty line.
\item If the command raises an error, the reply is a single line beginning with
      \verb+error:+ followed by the error message.
\end{itemize}

Two details matter when writing a client.  Values are rendered for the wire
rather than as Python would print them: a boolean is {\tt true} or
{\tt false} rather than {\tt True} or {\tt False}, and a string is returned
unquoted.  And only the value of the command reaches the socket --- anything
the command prints, and the traceback of an error, go to FEMM's own output
instead, so a client cannot recover them from the reply.

All of the supplied toolboxes are thin wrappers over this protocol, each
presenting it as native function calls in its host language: {\tt OctaveFEMM}
(MATLAB and Octave), {\tt MathFEMM} (Mathematica), {\tt pyfemm} (Python), and
equivalent packages for Julia, Scilab, and Lua.  Those packages serve both as
ready-to-use interfaces and as worked examples of how to speak the protocol
from a new client.
